\documentclass[tikz]{standalone}

\usepackage{tikz}
\usepackage{tkz-euclide}
\usepackage{fourier-otf}
\usepackage{fontspec}

\usetikzlibrary{calc}
\usetikzlibrary{3d}

\tikzset{
  every picture/.append style={scale=2, every node/.style={scale=2}}
}



\begin{document}
  \begin{tikzpicture}[scale=0.375]
    % Définition des limites
    \def\xmin{-9}
    \def\xmax{12}
    \def\ymin{-5}
    \def\ymax{8}
    
    % Axes
    \draw[->] (\xmin,0) -- (\xmax,0);
    \draw[->] (0,\ymin) -- (0,\ymax);
    
    % Graduations sur les axes
    \foreach \x in {\xmin,...,\xmax} {
        \ifnum\x\neq0
            \draw (\x,0) -- (\x,0.1) node[below, font=\tiny] {};
        \fi
    }
    \foreach \y in {\ymin,...,\ymax} {
        \ifnum\y\neq0
            \draw (0,\y) -- (0.1,\y) node[left, font=\tiny] {};
        \fi
    }
    
    % Origine
    \node[below left] at (0,0) {$O$};
    
    % Vecteurs unitaires
    \draw[thick, ->] (0,0) -- (1,0);
    \node[below] at (0.5,0) {\small $\vec{\imath}$};
    \draw[thick, ->] (0,0) -- (0,1);
    \node[left] at (0,0.5) {\small $\vec{\jmath}$};
    
    % Asymptotes
    \draw[blue] (-11,-11) -- (9,9); % asymptote oblique y=x
    \draw[blue] (1,\ymin) -- (1,\ymax); % asymptote verticale x=1
    
    % Étiquettes des asymptotes
    \node[right, rotate=90] at (1,-3) {$x=1$};
    \node[right, rotate=45] at (6,6) {$y=x$};
    
    % Fonction f(x) = x + 2/(x-1)
    % Partie gauche de la fonction (x < 1) - tracée deux fois comme dans l'original
    \draw[red, thick, domain=\xmin:-0.0001, samples=200] plot (\x, {\x + 2/(\x-1)});
    \draw[red, thick, domain=\xmin:\xmax, samples=200] plot (\x, {\x + 2/(\x-1)});
    
    % Note: La fonction est tracée en continu sauf à la discontinuité x=1
    
\end{tikzpicture}
\end{document}
